%%
%% Hypothesentest
%% 2023 - 05 - 24 φ
%%


\subsection{Fehler 1. Art / Fehler 2. Art}


\begin{definition}{Fehler 1. Art}{}\index{Fehler
    1. Art}\index{beta@$\beta$-Fehler}
  Wird die Nullhypothese $H_0$ fälschlicherweise verworfen, obschon
  sie wahr ist, so sprechen wir von einem $\alpha$-Fehler oder \textbf{Fehler erster Art}.
\end{definition}

\begin{bemerkung}{Fehler 1. Art}{}
  Fehler 1. Art bedeutet:

  Wir beschuldigen einen «fairen» Anbieter/Würfel/....
  \end{bemerkung}

Wir können für unser Pharma-Unternehmen, das weiterhin behauptet, dass
80\,\% der Patienten mit ihrem Medikament geheilt werden.

Unser Annahmebereich sei von 73 bis 100 Personen mit Linderung
festgesetz. Der Ablehnungsbereich somit von 0 bis 72 Personen.

Mit welcher Wahrscheinlichkeit wird in diesem Fall die $H_0$-Hypothese
fälschlicherweise verworfen, obschon das Medikament zu 80\,\% wirkt?

Mit welcher Wahrscheinlichkeit also landen wir im Ablehnungsbereich,
obschon das Medikament zu 80\,\% wirkt?

\TNTeop{
Selbst bei einer Wirkung von 80
  Von 0 bis max 72 werden  3.42\,\% aller Tests mit
  100 Probanden abschneiden. In diesen 3.42\,\% sprechen wir von einem
  Fehler 1. Art ($\alpha$-Fehler)

  Diese 3.42\% sind in unserem Fall gleichzeitig das Signifikanzniveau $\alpha$.

}%% end TNTeop (implicit \newpage)

\begin{definition}{Fehler 2. Art}{}\index{Fehler
    1. Art}\index{beta@$\beta$-Fehler}
  Wird die Nullhypothese $H_0$ fälschlicherweise angenommen (obschon
  $H_1$ wahr ist), so
  sprechen wir von einem $\beta$-Fehler oder \textbf{Fehler zweiter Art}.
\end{definition}
\begin{bemerkung}{Fehler 2. Art}{}
  Fehler 2. Art bedeutet:

  Wir lassen einen «Schuldigen», «gezinkten Würfel», ... laufen.

  Wir übersehen die Unwirksamkeit, den gezinkten Würfeln, ...
\end{bemerkung}


Nehmen wir diesmal wieder den Annahmebereich für unser Pharma-Test an 100
Probanden von 73 bis 100 Personen mit eintretenden Linderung.

Es hat sich aber herausgestellt, dass das Medikament nur mit 75\,\%
Linderung anbieten kann (die ersten Behauptungen mit 80\,\% des
Unternehmens waren falsch).

Mit welcher Wahrscheinlichkeit landen wir also fälschlicherweise im Annahmebereich (73-100),
obschon das Medikament nur zu 75\,\% wirkt?

\begin{bemerkung}{$\beta$-Fehler}{}
  Den Fehler 2 Art ($\beta$-Fehler) können wir erst rechnen,
  wenn wir die effektive, erst verborgene, Prozentzahl kennen.
\end{bemerkung}


\TNTeop{

  $$P(X \ge 73) = 1-P(X \le 72) \approx 1-27.76\,\% \approx 72.24\,\%$$

Wir können in 72.24\,\% aller Fälle nicht entscheiden, ob das
Medikament nun zu 80\,\% wirkt oder weniger.


}%% end TNTeop (implicit \newpage)
